\section{Governors and Control Loop}

\subsection{Loop Overview}

The control loop is single-threaded from the core perspective. The platform calls one step function at a fixed period.

\begin{lstlisting}[language=C,caption={Conceptual control-loop order}]
thermal_core_step(ctx):
    read sensors through platform snapshot
    update filters and zone temperatures
    evaluate trips
    run each zone governor
    arbitrate shared actuators
    apply policy modifiers
    apply slew limits
    emit telemetry and fault events
\end{lstlisting}

Figure~\ref{fig:closed-loop-control} shows the same flow as a closed-loop control system. The protocol-agnostic core owns the middle of the loop: filtering, zone state, trip evaluation, governor output, arbitration, policy modifiers, slew limits, faults, and telemetry. Platform-specific code only supplies snapshots and applies actuator requests.

\begin{figure}[htbp]
\centering
\resizebox{\textwidth}{!}{%
\begin{tikzpicture}[
    font=\small,
    node distance=0.72cm,
    block/.style={draw, rounded corners=2pt, align=center, minimum height=0.85cm, minimum width=2.35cm, fill=blue!6},
    coreblock/.style={block, fill=orange!10},
    ioblock/.style={block, fill=gray!12},
    plantblock/.style={block, fill=green!10, minimum width=2.8cm},
    contextblock/.style={block, fill=purple!8, minimum width=2.6cm},
    sum/.style={circle, draw, align=center, inner sep=1pt, minimum size=0.6cm, fill=white},
    arrow/.style={->, thick}
]
    \node[coreblock] (target) {Targets\\trips, setpoints};
    \node[sum, right=0.65cm of target] (compare) {$\Sigma$};
    \node[coreblock, right=0.65cm of compare] (governor) {Governor\\step-wise or PID};
    \node[coreblock, right=0.65cm of governor] (policy) {Arbitration + policy\\caps, offsets, slew};
    \node[ioblock, right=0.65cm of policy] (actuator) {Actuator adapter\\PWM / fan request};
    \node[plantblock, right=0.65cm of actuator] (plant) {Thermal plant\\device + cooling};

    \node[contextblock, above=0.75cm of policy] (context) {Context snapshot\\vehicle speed, mode};
    \node[contextblock, above=0.75cm of plant] (disturbance) {Disturbances\\ambient, workload};
    \node[ioblock, below=1.0cm of plant] (sensors) {Sensor snapshot\\temperature, tach};
    \node[coreblock, below=1.0cm of governor] (filters) {Filters + zone state\\temperature, faults};

    \draw[arrow] (target) -- (compare);
    \draw[arrow] (compare) -- (governor);
    \draw[arrow] (governor) -- (policy);
    \draw[arrow] (policy) -- (actuator);
    \draw[arrow] (actuator) -- (plant);
    \draw[arrow] (disturbance) -- (plant);
    \draw[arrow] (context) -- (policy);
    \draw[arrow] (plant) -- (sensors);
    \draw[arrow] (sensors.west) -- (filters.east);
    \draw[arrow] (filters.west) -| (compare.south);
\end{tikzpicture}%
}
\caption{Protocol-agnostic closed-loop thermal-control model}
\label{fig:closed-loop-control}
\end{figure}

\subsection{Governors}

The step-wise governor maps trip state to discrete cooling states. The PID governor targets a setpoint using Q16.16 fixed-point math by default. Anti-windup behavior and output clamping must be specified before final release.

\subsection{Arbitration}

When multiple zones request the same actuator, v1 uses max-wins arbitration. This is conservative and easy to reason about for a single shared fan.

\subsection{To Fill In Later}

\begin{itemize}[leftmargin=*]
    \item final pseudocode matching implementation order;
    \item fixed-point PID equations;
    \item anti-windup rule;
    \item unit-test references for step-wise and PID behavior.
\end{itemize}

\endinput
